% !TEX root = ../main.tex

\chapter{绪论}
\chaptereng{Introduction}
\section{研究背景与问题提出}

为满足国家2030年“碳达峰”、2060年“碳中和”的重大战略目标，加快构建清洁低碳、安全高效的能源体系已成为我国能源领域发展的核心方向。燃气轮机作为一种将化学能转化为机械能的热力旋转机械，凭借其功率密度大、启动快、效率高、污染低等显著优势，在船舶动力、分布式发电、天然气长输管线增压、航空推进以及工业驱动等国民经济关键领域发挥着不可替代的重要作用。

燃气轮机的基本工作原理可以概括为三个核心过程：压气机吸入并压缩空气，燃烧室中燃料与高压空气混合燃烧释放热能，高温燃气在涡轮中膨胀做功输出轴功率。在整个循环中，燃烧室是决定整机性能、排放特性和运行安全性的核心部件。当代燃气轮机正朝着更大功率、更高效率和更低排放的方向发展，压气机压比和涡轮前温度持续攀升，对燃烧室设计提出了更为严苛的要求。

为解决提高热效率和降低污染物排放之间的矛盾，轴向分级燃烧技术在先进燃气轮机中的应用越来越广泛。该技术将部分燃料和空气分流至二级燃烧区，以燃烧室壁面横向射流的形式加入预混燃料，在高温低氧环境中自燃着火，进一步提高燃烧室出口温度并进入透平做功。然而，二级横向射流火焰的稳定性是轴向分级燃烧面临的核心挑战之一——火焰吹熄将直接导致燃烧效率下降、排放恶化，严重时可能造成燃烧室熄火停车。

在先进燃气轮机中，燃烧室下游与涡轮之间的过渡流道普遍采用收缩设计，用于优化流场过渡、提高涡轮进口参数。这一几何特征可能通过改变主流速度、压力梯度和剪切层强度影响上游横向射流火焰的稳定特性。然而，现有关于横向射流火焰熄火的研究大多在直通道或恒定截面横流环境中进行，对下游收缩流道在热流工况下的定量影响及其机理关注不足。

针对上述工程问题，本文着重围绕下游收缩流道对横向射流火焰熄火边界的定量影响及不同收缩高度的调控趋势进行研究。

\section{研究方法与技术路线}

本研究采用热态实验与数值模拟相结合的方法，以热态实验为主、数值模拟为辅。实验依托课题组自研轴向分级燃烧器平台，通过一级预混燃烧构建稳定高温主流；二级燃烧室出口安装可更换收缩体，实现下游流道几何参数（收缩型线、收缩比）的精准调控。结合高速摄影、紫外火焰探测器、热电偶和质量流量控制系统，同步采集火焰形态、火焰信号、温度场及流量参数。数值模拟采用Cantera软件构建简化反应流模型，预估热流工况范围并辅助解释实验现象。

技术路线为：搭建轴向分级燃烧器平台→设计加工收缩体→建立同温工况筛选方法→开展对比实验→建立占空比判据→分析数据→揭示影响规律。整个研究流程遵循“先跑通、再优化”的思路，首先确保实验系统能够稳定运行并获取初步数据，在此基础上逐步完善测量方法和实验方案。

关键技术问题包括：火焰探测器高频采集（解决商用探测器信号平滑化问题）、同温工况筛选（解耦主流温度和速度的影响）、占空比判据建立（将火焰稳定性从定性观察转化为定量指标）

\section{主要研究内容与论文结构}

本文共分为七章，各章内容安排如下：

第1章为引言，介绍研究背景，提出本文围绕吹熄测量方法建立、收缩流道定量影响、测量方法改进探索三个核心问题展开，给出全文结构导航。

第2章为文献综述，系统回顾横向射流冷态混合、火焰稳定机制、吹熄理论与模型、流道结构对火焰稳定性影响四方面的研究进展，指出现有研究的不足与本课题的定位。

第3章为实验平台与测量方法，详细描述轴向分级燃烧器的各组成部分、收缩体的设计加工过程、测量系统的构成与改进。

第4章为测量方法的改进探索，针对商用4-20mA探测器响应速度慢、丢失高频信息的局限性，提出基于日盲紫外光电管C7027A的高频脉冲信号采集方案，涵盖电路设计、示波器验证及Python脉冲检测算法开发。

第5章热流工况下横向射流火焰吹熄测量方法的建立，包括同温工况筛选、二级射流燃料优化、火焰探测器信号采集改进、占空比判据的提出与验证。

第6章下游收缩流道对熄火边界的定量影响，对比无收缩段与60mm、32mm两种收缩段构型的火焰稳定性差异，初步探讨收缩高度对火焰稳定性的调控趋势。

第7章为结论与展望，汇总本研究的阶段性结论，总结创新点，提出后续工作方向。